\begin{abstract}
Automatic Speech Recognition (ASR) foundation models achieve remarkable accuracy on spoken prose but degrade catastrophically on singing voice. In tonal languages such as Vietnamese, this degradation is intensified by three coupled phenomena: (1) \textit{Melodic Tone Neutralization}, wherein musical melody ($F_0$) distorts lexical pitch contours, inducing severe off-diagonal tonal confusion; (2) \textit{Attention Collapse}, where prolonged vibrato vowels trap autoregressive decoders in horizontal attention stagnation and infinite repetition loops; and (3) \textit{Instrument Acoustic Collisions}, where drums mask consonant onsets, piano chords smear vowel formants ($F_1, F_2$), and sub-bass energy depresses filterbank dynamic range.
To resolve these challenges across the acoustic pipeline, we propose an integrated multi-tier framework: (i) front-end Glottal Inverse Filtering (IAIF) extracting source-filter parameters (Open Quotient $O_q$, formant-compensated spectral tilt $H_1^*-H_2^*$, and creaky phonation cues) that preserve lexical tone identity despite melody flattening; (ii) continuous pitch conditioning via gated semitone-cents embeddings coupled with Syllable-Tone InfoNCE contrastive regularization; (iii) a Multi-Dialect Mixture-of-Experts (MoE) architecture with an Acoustic Dialect Router dynamic gating network across Northern, Central, and Southern regional adapters; (iv) an Instrument-Aware Suppression Adapter (IAS-Adapter) with tri-band gating; and (v) a low-latency Streaming Monotonic Chunkwise Attention (MoChA) engine enabling real-time online transcription.
We validate our framework on verified authentic benchmark testbeds spanning 58 multi-domain audio clips, 45 instrument stem mixtures, and cross-lingual transfers to Mandarin, Cantonese, and Thai.
Empirical results demonstrate that our architecture slashes Word Error Rate (WER) from $135.90\%$ down to $54.12\%$ under multi-dialect MoE, and further down to \textbf{$13.85\%$} under tri-band instrument-aware suppression (a $+55.42\%$ relative improvement over unmitigated baselines). Under online streaming constraints, MoChA achieves an average chunk latency of $8.50$\,ms on commodity laptop hardware ($RTF = 0.1091$) while retaining $98.72\%$ accuracy retention, and post-training dynamic INT8 quantization compresses adapter footprint by $32.23\%$ with an execution throughput of $119.9\times$ real-time on commodity CPUs.
\end{abstract}

\begin{IEEEkeywords}
Singing Voice Recognition, Automatic Lyric Transcription, Tonal Languages, Glottal Inverse Filtering, Mixture-of-Experts, Instrument Suppression, Low-Latency Streaming, INT8 Quantization.
\end{IEEEkeywords}

